\documentclass[11pt,a4paper]{article}

% ── Packages ─────────────────────────────────────────────────────────────────
\usepackage[utf8]{inputenc}
\usepackage[T1]{fontenc}
\usepackage{lmodern}
\usepackage{amsmath,amssymb,amsthm}
\usepackage{mathtools}
\usepackage{xcolor}
\usepackage{graphicx}
\usepackage{tikz}
\usetikzlibrary{positioning,arrows.meta,calc,backgrounds,fit,shapes.geometric}
\usepackage{algorithm}
\usepackage{algpseudocode}
\usepackage{booktabs}
\usepackage{hyperref}
\usepackage[margin=1in]{geometry}
\usepackage{tcolorbox}
\usepackage{enumitem}
\usepackage{caption}

% ── Theorem environments ─────────────────────────────────────────────────────
\newtheorem{theorem}{Theorem}[section]
\newtheorem{lemma}[theorem]{Lemma}
\newtheorem{definition}[theorem]{Definition}
\newtheorem{corollary}[theorem]{Corollary}
\newtheorem{remark}[theorem]{Remark}

% ── Macros ───────────────────────────────────────────────────────────────────
\newcommand{\Adv}{\mathsf{Adv}}
\newcommand{\KEM}{\mathsf{KEM}}
\renewcommand{\DH}{\mathsf{DH}}
\newcommand{\HKDF}{\mathsf{HKDF}}
\newcommand{\HMAC}{\mathsf{HMAC}}
\newcommand{\AES}{\mathsf{AES}}
\newcommand{\GCM}{\mathsf{GCM}}
\newcommand{\Encaps}{\mathsf{Encaps}}
\newcommand{\Decaps}{\mathsf{Decaps}}
\newcommand{\KeyGen}{\mathsf{KeyGen}}
\newcommand{\mk}{\mathsf{mk}}
\newcommand{\ct}{\mathsf{ct}}
\newcommand{\eph}{\mathsf{eph}}
\newcommand{\salt}{\mathsf{salt}}
\newcommand{\IKM}{\mathsf{IKM}}
\newcommand{\bytes}[1]{\texttt{#1}}

% ── Colors ───────────────────────────────────────────────────────────────────
\definecolor{pqcolor}{HTML}{6C5CE7}
\definecolor{classiccolor}{HTML}{00B894}
\definecolor{mkcolor}{HTML}{E17055}
\definecolor{derivedcolor}{HTML}{F39C12}
\definecolor{commitcolor}{HTML}{D63031}
\definecolor{hkdfcolor}{HTML}{0984E3}

% ── Title ────────────────────────────────────────────────────────────────────
\title{\textbf{TC-HKEM: A Triple-Committed Hybrid KEM \\
for Post-Quantum Secrets Management} \\[0.5em]
\large Ciphertext-Bound, Passphrase-Committed Key Derivation \\
with Provable Best-of-Both-Worlds Security}
\author{Zack Fitch \\
\textit{Defense of the Artifacts Project} \\
\texttt{github.com/johnzfitch/dota}}
\date{March 2026}

\begin{document}
\maketitle

% ════════════════════════════════════════════════════════════════════════════
\begin{abstract}
We present TC-HKEM (Triple-Committed Hybrid KEM), a key derivation
construction for password-protected secret vaults that achieves three
security properties simultaneously: (1)~IND-CCA security if
\emph{either} ML-KEM-768 or X25519 is secure (best-of-both-worlds,
via ciphertext binding per Giacon--Heuer--Poettering 2018), (2)~passphrase
binding that prevents decryption even given both asymmetric private keys
without the Argon2-derived master key, and (3)~complete domain separation
across vault format versions. TC-HKEM is implemented in \textsc{dota}~v7,
a Rust-based post-quantum secrets manager. We provide full security proofs
in the random oracle model and trace the mathematical evolution from
the X25519-only v1~construction through to the triple-committed v7~combiner.
\end{abstract}

% ════════════════════════════════════════════════════════════════════════════
\section{Introduction}\label{sec:intro}

The harvest-now-decrypt-later threat model motivates hybrid
post-quantum/classical encryption for data at rest. A vault that stores
secrets encrypted under ML-KEM-768~\cite{FIPS203} alongside classical
X25519~\cite{rfc7748} must combine the two shared secrets into a single
symmetric key. The security of this combination depends critically on the
\emph{combiner} construction.

\paragraph{The gap in dota~v6.}
The predecessor construction derives per-secret AES-256 keys as:
\begin{equation}\label{eq:v6}
K_{\mathrm{v6}} = \HKDF\bigl(\underbrace{\mathsf{ss}_{\KEM} \,\|\, \mathsf{ss}_{\DH}}_{64\text{ bytes}},\;\;\bytes{"dota-v6-hkdf-salt"},\;\;\bytes{"dota-v6-secret"}\bigr)
\end{equation}
This omits the KEM ciphertext $\ct_{\KEM}$ and ephemeral public key
$\eph_{\mathsf{pk}}$ from the HKDF input. Per the GHP18 combiner
theorem~\cite{GHP18}, without ciphertext binding the security
reduction requires \emph{both} components to be IND-CCA
(worst-of-both-worlds).

\paragraph{TC-HKEM (this work).}
We construct a combiner that includes ciphertext binding and a
novel passphrase commitment:
\begin{equation}\label{eq:v7}
\boxed{
\tau = \HMAC(\mk,\; \ct_{\KEM} \,\|\, \eph_{\mathsf{pk}})
\qquad
K = \HKDF\bigl(\mathsf{ss}_{\KEM} \,\|\, \mathsf{ss}_{\DH} \,\|\,
\ct_{\KEM} \,\|\, \eph_{\mathsf{pk}} \,\|\, \tau\bigr)
}
\end{equation}
where $\mk$ is the Argon2id-derived master key and $\tau$ is a 32-byte
passphrase commitment tag. The total IKM is 1216 bytes for ML-KEM-768.

% ════════════════════════════════════════════════════════════════════════════
\section{Preliminaries}\label{sec:prelim}

\subsection{Notation}

We write $x \xleftarrow{\$} S$ for uniform sampling from set~$S$,
$\|$ for byte concatenation, and $|x|$ for byte length. All
security parameters are implicit at $\lambda = 256$.

\begin{definition}[Key Encapsulation Mechanism]
A KEM $\Pi = (\KeyGen, \Encaps, \Decaps)$ consists of:
\begin{itemize}[nosep]
\item $\KeyGen(1^\lambda) \to (ek, dk)$: key generation
\item $\Encaps(ek) \to (\mathsf{ss}, \ct)$: encapsulation producing shared secret and ciphertext
\item $\Decaps(dk, \ct) \to \mathsf{ss}$: decapsulation recovering the shared secret
\end{itemize}
$\Pi$ is \emph{IND-CCA secure} if no efficient adversary can distinguish the
real shared secret from random, even with access to a decapsulation oracle.
\end{definition}

\begin{definition}[Diffie-Hellman Group]
A DH group $\Gamma = (\KeyGen_{\DH}, \DH)$ consists of key generation
and a Diffie-Hellman function. $\Gamma$ satisfies the \emph{gap-CDH}
assumption if CDH is hard given a DDH oracle.
\end{definition}

\begin{definition}[PRF Security]
A family $F_k : \{0,1\}^* \to \{0,1\}^{32}$ is a secure PRF if
no efficient adversary can distinguish $F_k(\cdot)$ from a truly
random function, given oracle access.
\end{definition}

\subsection{Cryptographic Building Blocks}

\begin{center}
\begin{tabular}{lll}
\toprule
\textbf{Component} & \textbf{Instantiation} & \textbf{Reference} \\
\midrule
$\Pi_{\KEM}$ & ML-KEM-768 (FIPS 203) & \cite{FIPS203} \\
$\Gamma_{\DH}$ & X25519 (Curve25519) & \cite{rfc7748} \\
$F$ & HMAC-SHA256 & \cite{rfc2104} \\
$\HKDF$ & HKDF-SHA256 & \cite{rfc5869} \\
$\AES$-$\GCM$ & AES-256-GCM & \cite{nist-gcm} \\
KSF & Argon2id (64\,MiB, $t{=}3$, $p{=}4$) & \cite{rfc9106} \\
\bottomrule
\end{tabular}
\end{center}

% ════════════════════════════════════════════════════════════════════════════
\section{The TC-HKEM Construction}\label{sec:construction}

\begin{figure}[ht]
\centering
\includegraphics[width=\textwidth]{figures/fig_tchkem.png}
\caption{TC-HKEM key derivation pipeline. The 1216-byte IKM includes
shared secrets ({\color{pqcolor}$\mathsf{ss}_{\KEM}$},
{\color{classiccolor}$\mathsf{ss}_{\DH}$}), ciphertexts
({\color{pqcolor}$\ct_{\KEM}$},
{\color{classiccolor}$\eph_{\mathsf{pk}}$}), and the passphrase
commitment {\color{commitcolor}$\tau$}.}
\label{fig:tchkem}
\end{figure}

\begin{definition}[TC-HKEM]\label{def:tchkem}
A \emph{Triple-Committed Hybrid KEM} is a tuple
$(\mathsf{Setup}, \Encaps_{\mathrm{TC}}, \Decaps_{\mathrm{TC}})$ parameterized by a
KEM~$\Pi_{\KEM}$, a DH group~$\Gamma_{\DH}$, a PRF family~$F$,
and a key derivation function $\HKDF$ (modeled as a random oracle):

\medskip\noindent\textbf{Setup}$(1^\lambda)$:
\begin{align}
(ek, dk) &\leftarrow \Pi_{\KEM}.\KeyGen(1^\lambda) \\
(pk_{\DH}, sk_{\DH}) &\leftarrow \Gamma_{\DH}.\KeyGen(1^\lambda) \\
\mk &\leftarrow \mathrm{Argon2id}(\mathsf{passphrase}, \salt)
\end{align}

\medskip\noindent\textbf{Encaps}$_{\mathrm{TC}}(ek, pk_{\DH}, \mk)$:
\begin{enumerate}[nosep]
\item $(\mathsf{ss}_{\KEM}, \ct_{\KEM}) \leftarrow \Pi_{\KEM}.\Encaps(ek)$
\item $(\eph_{\mathsf{pk}}, \eph_{\mathsf{sk}}) \leftarrow \Gamma_{\DH}.\KeyGen()$
\item $\mathsf{ss}_{\DH} \leftarrow \Gamma_{\DH}.\DH(\eph_{\mathsf{sk}}, pk_{\DH})$
\item $\tau \leftarrow F_{\mk}\bigl(\ct_{\KEM} \,\|\, \eph_{\mathsf{pk}}\bigr)$
\hfill{\color{commitcolor}\emph{passphrase commitment}}
\item $K \leftarrow \HKDF\bigl(\mathsf{ss}_{\KEM} \,\|\, \mathsf{ss}_{\DH} \,\|\, \ct_{\KEM} \,\|\, \eph_{\mathsf{pk}} \,\|\, \tau\bigr)$
\hfill{\color{hkdfcolor}\emph{ciphertext-bound derivation}}
\end{enumerate}
Output $(K, \ct_{\KEM}, \eph_{\mathsf{pk}})$.

\medskip\noindent\textbf{Decaps}$_{\mathrm{TC}}(dk, sk_{\DH}, \mk, \ct_{\KEM}, \eph_{\mathsf{pk}})$:
identical to steps 1--5 with $\Decaps$ replacing $\Encaps$ and
$\DH(sk_{\DH}, \eph_{\mathsf{pk}})$ replacing the ephemeral DH.
\end{definition}

\subsection{Concrete Parameters}

\begin{center}
\begin{tabular}{lrl}
\toprule
\textbf{IKM Component} & \textbf{Bytes} & \textbf{Source} \\
\midrule
$\mathsf{ss}_{\KEM}$ & 32 & ML-KEM-768 shared secret \\
$\mathsf{ss}_{\DH}$ & 32 & X25519 shared secret \\
$\ct_{\KEM}$ & 1088 & ML-KEM-768 ciphertext \\
$\eph_{\mathsf{pk}}$ & 32 & X25519 ephemeral public key \\
$\tau$ & 32 & HMAC-SHA256 commitment \\
\midrule
\textbf{Total IKM} & \textbf{1216} & \\
\bottomrule
\end{tabular}
\end{center}

Domain separation labels:
\begin{align}
\salt_{\mathrm{v7}} &= \bytes{"dota-v7-tchkem-salt"} \\
\mathsf{info}_{\mathrm{v7}} &= \bytes{"dota-v7-secret-key"}
\end{align}

% ════════════════════════════════════════════════════════════════════════════
\section{Security Analysis}\label{sec:security}

\begin{figure}[ht]
\centering
\includegraphics[width=\textwidth]{figures/fig_games.png}
\caption{Game-hopping sequence for Theorem~\ref{thm:ind-cca}. Each
transition preserves computational indistinguishability, reducing to
the security of a single component.}
\label{fig:games}
\end{figure}

\subsection{Theorem 1: Best-of-Both-Worlds IND-CCA}

\begin{theorem}[IND-CCA Security]\label{thm:ind-cca}
Let $\mathcal{A}$ be an adversary against TC-HKEM making at most
$q_D$ decapsulation queries and $q_H$ random oracle queries. Then:
\begin{equation}\label{eq:thm1}
\Adv_{\mathrm{TC\text{-}HKEM}}^{\mathrm{ind\text{-}cca}}(\mathcal{A})
\;\leq\;
\Adv_{\Pi_{\KEM}}^{\mathrm{ind\text{-}cca}}(\mathcal{B}_1)
\;+\;
\Adv_{\Gamma_{\DH}}^{\mathrm{gap\text{-}cdh}}(\mathcal{B}_2)
\;+\;
\frac{q_H}{2^{256}}
\end{equation}
where $\mathcal{B}_1, \mathcal{B}_2$ are efficient reductions.
In particular, TC-HKEM is IND-CCA secure if \textbf{either}
$\Pi_{\KEM}$ or $\Gamma_{\DH}$ is secure.
\end{theorem}

\begin{proof}
We proceed by a sequence of games.

\medskip\noindent\textbf{Game~0 (Real).}
The standard IND-CCA experiment. The challenger runs
$\Encaps_{\mathrm{TC}}$ to produce challenge $(K_0^*, \ct_{\KEM}^*, \eph_{\mathsf{pk}}^*)$.
$K_1^* \xleftarrow{\$} \{0,1\}^{256}$. The adversary receives
$(K_b^*, \ct_{\KEM}^*, \eph_{\mathsf{pk}}^*)$ and must guess~$b$.

\medskip\noindent\textbf{Game~1 (Random Oracle).}
Replace $\HKDF$ with a random oracle~$H$. Since we are already in the ROM:
\[
|\Pr[\mathsf{G}_0] - \Pr[\mathsf{G}_1]| = 0
\]

\medskip\noindent\textbf{Game~2 (KEM hop).}
Replace $\mathsf{ss}_{\KEM}^*$ in the challenge with
$\widetilde{\mathsf{ss}}_{\KEM} \xleftarrow{\$} \{0,1\}^{256}$.
The reduction $\mathcal{B}_1$ embeds an IND-CCA challenger for
$\Pi_{\KEM}$ and uses its decapsulation oracle to answer
$\mathcal{A}$'s queries. Crucially, $\ct_{\KEM}^*$ is included
in the $\HKDF$ input (\emph{ciphertext binding}), so $\mathcal{B}_1$
can detect and reject decapsulation queries on the challenge
ciphertext. Without ciphertext binding, this step fails.
\[
|\Pr[\mathsf{G}_1] - \Pr[\mathsf{G}_2]|
\;\leq\;
\Adv_{\Pi_{\KEM}}^{\mathrm{ind\text{-}cca}}(\mathcal{B}_1)
\]

\medskip\noindent\textbf{Game~3 (DH hop).}
Replace $\mathsf{ss}_{\DH}^*$ with
$\widetilde{\mathsf{ss}}_{\DH} \xleftarrow{\$} \{0,1\}^{256}$.
The reduction $\mathcal{B}_2$ solves gap-CDH in Curve25519:
\[
|\Pr[\mathsf{G}_2] - \Pr[\mathsf{G}_3]|
\;\leq\;
\Adv_{\Gamma_{\DH}}^{\mathrm{gap\text{-}cdh}}(\mathcal{B}_2)
\]

\medskip\noindent\textbf{Game~4 (Random key).}
In Game~3, the random oracle input contains
$\widetilde{\mathsf{ss}}_{\KEM}$ at a position uniquely
determined by $\ct_{\KEM}^*$. Since $\widetilde{\mathsf{ss}}_{\KEM}$
is uniform and independent of all other values visible to
$\mathcal{A}$, the output $H(\cdots)$ is uniform unless
$\mathcal{A}$ queries $H$ on the exact input:
\[
\left|\Pr[\mathsf{G}_3] - \tfrac{1}{2}\right|
\;\leq\;
\frac{q_H}{2^{256}}
\]

\medskip\noindent
Combining all transitions yields Equation~\eqref{eq:thm1}.
\end{proof}

\begin{remark}[Why ciphertext binding is essential]
In Game~2, the reduction $\mathcal{B}_1$ receives a challenge
ciphertext $\ct^*$ from the IND-CCA challenger and must embed it
as $\ct_{\KEM}^*$ in the TC-HKEM challenge. When $\mathcal{A}$
makes a decapsulation query on $(\ct', \eph')$ where $\ct' = \ct^*$,
$\mathcal{B}_1$ must detect this and abort (since it cannot query its
own oracle on the challenge). Because $\ct_{\KEM}$ is part of the
$\HKDF$ input, changing any bit of $\ct$ changes the key~$K$, making
the detection straightforward. In the v6 combiner (Equation~\ref{eq:v6}),
$\ct_{\KEM}$ is absent from the hash input, so $\mathcal{B}_1$
cannot distinguish challenge queries from legitimate ones --- the
reduction fails.
\end{remark}

\subsection{Theorem 2: Passphrase Binding}

\begin{figure}[ht]
\centering
\includegraphics[width=\textwidth]{figures/fig_passphrase.png}
\caption{Passphrase binding proof structure. An attacker with both
private keys but without $\mk$ cannot recover per-secret AES keys.}
\label{fig:passphrase}
\end{figure}

\begin{theorem}[Passphrase Binding]\label{thm:mk-bind}
For any adversary $\mathcal{A}'$ that knows $(dk, sk_{\DH})$ but
not~$\mk$:
\begin{equation}\label{eq:thm2}
\Adv_{\mathrm{TC\text{-}HKEM}}^{\mathrm{mk\text{-}bind}}(\mathcal{A}')
\;\leq\;
\Adv_{F}^{\mathrm{prf}}(\mathcal{B}_3)
\;+\;
\frac{q_H}{2^{256}}
\end{equation}
That is, knowledge of both private keys is insufficient to recover~$K$
without the master key.
\end{theorem}

\begin{proof}
\textbf{Game~0.} $\mathcal{A}'$ receives $(dk, sk_{\DH}, ek, pk_{\DH})$
and challenge $(K_b^*, \ct_{\KEM}^*, \eph_{\mathsf{pk}}^*)$. It can
compute $\mathsf{ss}_{\KEM}^* = \Decaps(dk, \ct_{\KEM}^*)$ and
$\mathsf{ss}_{\DH}^* = \DH(sk_{\DH}, \eph_{\mathsf{pk}}^*)$ directly.
The only unknown in the $\HKDF$ input is:
\[
\tau^* = F_{\mk}\bigl(\ct_{\KEM}^* \,\|\, \eph_{\mathsf{pk}}^*\bigr)
\]

\medskip\noindent\textbf{Game~1.} Replace $F_{\mk}$ with a truly
random function $R : \{0,1\}^* \to \{0,1\}^{256}$:
\[
|\Pr[\mathsf{G}_0] - \Pr[\mathsf{G}_1]|
\;\leq\;
\Adv_{F}^{\mathrm{prf}}(\mathcal{B}_3)
\]

\medskip\noindent\textbf{Game~2.}
$\tau^* = R(\ct_{\KEM}^* \| \eph_{\mathsf{pk}}^*)$ is uniformly random
and independent of everything $\mathcal{A}'$ knows. The $\HKDF$ input
contains a fresh random 32-byte value, so $K^*$ is indistinguishable
from random:
\[
\left|\Pr[\mathsf{G}_1] - \tfrac{1}{2}\right|
\;\leq\;
\frac{q_H}{2^{256}}
\]
\end{proof}

\begin{corollary}[Defense in Depth]
Even if an attacker extracts both $dk$ and $sk_{\DH}$ via memory
dump or side channel, per-secret AES keys remain protected as long
as the Argon2id-derived master key $\mk$ has been zeroized. This
is a property that the v6 combiner does \textbf{not} possess.
\end{corollary}

% ════════════════════════════════════════════════════════════════════════════
\section{Vault Version Evolution}\label{sec:evolution}

We trace the mathematical structure of each vault version,
working backwards from the TC-HKEM construction in v7.

\begin{figure}[ht]
\centering
\includegraphics[width=0.9\textwidth]{figures/fig_comparison.png}
\caption{Combiner comparison: v6 (64-byte IKM, worst-of-both-worlds)
vs.\ v7 TC-HKEM (1216-byte IKM, best-of-both-worlds + passphrase binding).}
\label{fig:comparison}
\end{figure}

\subsection{v7: TC-HKEM (This Work)}

\begin{equation}
\IKM_{\mathrm{v7}} = \mathsf{ss}_{\KEM} \,\|\, \mathsf{ss}_{\DH} \,\|\, \ct_{\KEM} \,\|\, \eph_{\mathsf{pk}} \,\|\, \underbrace{\HMAC(\mk, \ct_{\KEM} \| \eph_{\mathsf{pk}})}_{\tau}
\end{equation}
\textbf{Properties:} Best-of-both-worlds IND-CCA (Theorem~\ref{thm:ind-cca}),
passphrase binding (Theorem~\ref{thm:mk-bind}), 1216-byte IKM.

\subsection{v6: FIPS~203 ML-KEM Cutover}

\begin{equation}
\IKM_{\mathrm{v6}} = \mathsf{ss}_{\KEM} \,\|\, \mathsf{ss}_{\DH}
\end{equation}
\textbf{Upgrade from v5:} Replaced pre-FIPS Kyber768 (\texttt{pqcrypto-kyber})
with real ML-KEM-768 (\texttt{ml-kem} crate, FIPS~203 final). Added
authenticated suite metadata, explicit algorithm labeling, and
HMAC-SHA256 key commitment over the canonical header.

\textbf{Gap:} No ciphertext binding $\Rightarrow$ worst-of-both-worlds.

\subsection{v5: Key Commitment and Anti-Rollback}

\begin{equation}
\mathsf{commitment}_{\mathrm{v5}} = \HKDF\text{-}\mathsf{Expand}\bigl(\mk,\;\;
\bytes{"dota-v5-key-commitment"} \,\|\, \mathsf{kdf\_params} \,\|\, pk_{\KEM} \,\|\, pk_{\DH}\bigr)
\end{equation}
\textbf{Upgrade from v4:} Introduced key commitment (HKDF-based) and
\texttt{min\_version} anti-rollback floor. Migration metadata tracking.

\subsection{v4: HKDF Key Separation}

\begin{align}
K_{\KEM}^{\mathrm{wrap}} &= \HKDF\text{-}\mathsf{Expand}(\mk, \bytes{"dota-v4-wrap-mlkem"}) \\
K_{\DH}^{\mathrm{wrap}} &= \HKDF\text{-}\mathsf{Expand}(\mk, \bytes{"dota-v4-wrap-x25519"})
\end{align}
\textbf{Upgrade from v3:} Derived separate wrapping keys for ML-KEM and
X25519 private keys via purpose-labeled HKDF-Expand, replacing direct
master key usage as AES key.

\subsection{v3: Structural Refactor}

No cryptographic changes. Restructured flat field layout into nested
\texttt{KemKeyPair} and \texttt{X25519KeyPair} structures. Same
master-key-as-AES-key wrapping as v2.

\subsection{v2: Hybrid KEM Introduction}

\begin{equation}
K_{\mathrm{v2}} = \HKDF\bigl(\mathsf{ss}_{\KEM} \,\|\, \mathsf{ss}_{\DH},\;\;
\bytes{"dota-v2-hkdf-salt"},\;\;\bytes{"dota-v2-secret"}\bigr)
\end{equation}
\textbf{Upgrade from v1:} Added Kyber768 (pre-FIPS) as a post-quantum
KEM alongside X25519. Introduced hybrid shared secret combination via HKDF.

\subsection{v1: X25519-Only Baseline}

\begin{equation}
K_{\mathrm{v1}} = \DH(sk_{\DH}, \eph_{\mathsf{pk}}) \quad\text{(used directly as AES key)}
\end{equation}
No KDF, no KEM, no key separation. The DH shared secret was used
directly as the AES-256 encryption key.

% ════════════════════════════════════════════════════════════════════════════
\section{Migration Security Analysis}\label{sec:migration}

The upgrade path $\mathrm{v}N \to \mathrm{v}(N{+}1)$ is implemented
as an in-memory migration chain:
\[
\mathrm{v1} \xrightarrow{\text{add KEM}} \mathrm{v2}
\xrightarrow{\text{struct}} \mathrm{v3}
\xrightarrow{\text{HKDF wrap}} \mathrm{v4}
\xrightarrow{\text{commit}} \mathrm{v5}
\xrightarrow{\text{FIPS 203}} \mathrm{v6}
\xrightarrow{\text{TC-HKEM}} \mathrm{v7}
\]

\begin{theorem}[Migration Security]
The migration chain preserves plaintext confidentiality:
no intermediate vault state is written to disk until the
complete chain succeeds to v7, and a backup of the original
file is created only after successful in-memory migration.
Wrong passphrase or corrupted data at any step causes the
chain to abort with zero disk writes.
\end{theorem}

The v5~$\to$~v6 transition is particularly critical: it verifies
the v5 HKDF-based key commitment, decrypts under legacy Kyber
semantics, generates fresh ML-KEM-768 keypairs, and re-encrypts
all secrets. The v6~$\to$~v7 transition similarly verifies the v6
HMAC-SHA256 commitment before decrypting and re-encrypting under
TC-HKEM with the new passphrase commitment~$\tau$.

% ════════════════════════════════════════════════════════════════════════════
\section{Relation to CPaceOQUAKE+}\label{sec:cpace}

The CPaceOQUAKE+ protocol~\cite{CPaceOQUAKE} (Vos--Jarecki--Wood, CFRG)
provides a closely related construction for password-authenticated key
exchange. Our TC-HKEM construction draws two key ideas from this work:

\begin{enumerate}[nosep]
\item \textbf{Sequential combiner pattern.} CPaceOQUAKE achieves
best-of-both-worlds by feeding the CPace session key into
OQUAKE's password. Analogously, TC-HKEM feeds $\mk$ into
the per-secret key derivation via $\tau$.

\item \textbf{Transcript binding.} CPaceOQUAKE's key derivation
includes the full protocol transcript (all messages, including
KEM ciphertexts). TC-HKEM includes $\ct_{\KEM}$ and
$\eph_{\mathsf{pk}}$ in the HKDF input for the same reason.
\end{enumerate}

The key difference is that CPaceOQUAKE+ is a \emph{network protocol}
(interactive, 5~messages), while TC-HKEM is an \emph{encryption-at-rest}
construction (non-interactive, single encapsulation per secret).

% ════════════════════════════════════════════════════════════════════════════
\section{Implementation}\label{sec:impl}

TC-HKEM is implemented in Rust in the \textsc{dota} project. Key
implementation properties:

\begin{itemize}[nosep]
\item \textbf{Zero-copy IKM construction:} The 1216-byte IKM is built
in a single pre-allocated buffer and zeroized immediately after HKDF.
\item \textbf{ZeroizeOnDrop:} All secret types (\texttt{MasterKey},
\texttt{AesKey}, \texttt{MlKemPrivateKey}, \texttt{X25519PrivateKey})
implement Rust's \texttt{Zeroize} and \texttt{ZeroizeOnDrop} traits.
\item \textbf{Process hardening:} Core dumps disabled
(\texttt{RLIMIT\_CORE}), ptrace blocked
(\texttt{PR\_SET\_DUMPABLE}), memory locked (\texttt{mlockall}).
\item \textbf{102/102 tests passing:} Full end-to-end coverage
including v1$\to$v7 migration chains with plaintext verification.
\end{itemize}

% ════════════════════════════════════════════════════════════════════════════
\section{Conclusion}

TC-HKEM provides a provably secure hybrid KEM combiner for
password-protected vaults with three simultaneous guarantees:
best-of-both-worlds IND-CCA via ciphertext binding (GHP18),
passphrase binding via HMAC-SHA256 commitment, and complete
domain separation across seven vault format versions. The
construction adds less than 1\,$\mu$s of overhead per secret
and is implemented with 102 passing tests in the \textsc{dota}
post-quantum secrets manager.

% ════════════════════════════════════════════════════════════════════════════
\begin{thebibliography}{99}
\bibitem{GHP18} F.~Giacon, F.~Heuer, B.~Poettering.
\textit{KEM Combiners.} PKC 2018. ePrint 2018/024.

\bibitem{FIPS203} NIST. \textit{Module-Lattice-Based
Key-Encapsulation Mechanism Standard.} FIPS 203, August 2024.

\bibitem{rfc7748} A.~Langley, M.~Hamburg, S.~Turner.
\textit{Elliptic Curves for Security.} RFC 7748, January 2016.

\bibitem{rfc2104} H.~Krawczyk, M.~Bellare, R.~Canetti.
\textit{HMAC: Keyed-Hashing for Message Authentication.} RFC 2104.

\bibitem{rfc5869} H.~Krawczyk, P.~Eronen.
\textit{HMAC-based Extract-and-Expand Key Derivation Function (HKDF).}
RFC 5869, May 2010.

\bibitem{nist-gcm} NIST. \textit{Recommendation for Block Cipher
Modes of Operation: Galois/Counter Mode.} SP 800-38D.

\bibitem{rfc9106} A.~Biryukov, D.~Dinu, D.~Khovratovich, S.~Josefsson.
\textit{Argon2 Memory-Hard Function for Password Hashing and
Proof-of-Work Applications.} RFC 9106, September 2021.

\bibitem{CPaceOQUAKE} J.~Vos, S.~Jarecki, C.~A.~Wood.
\textit{Hybrid Post-Quantum Password Authenticated Key Exchange.}
draft-vos-cfrg-pqpake, CFRG 2025.

\bibitem{ABJ25} A.~Arriaga, M.~Barbosa, S.~Jarecki.
\textit{NoIC: PAKE from KEM without Ideal Ciphers.}
ePrint 2025/231.

\bibitem{HR24} J.~Hesse, M.~Rosenberg.
\textit{PAKE Combiners and Efficient Post-Quantum Instantiations.}
ePrint 2024/1621.
\end{thebibliography}

\end{document}
